\documentclass{article}
\usepackage{amsmath,amssymb}
\begin{document}

\section*{Higham Chapter 4 Kahan returned-sum coefficient blocker}

I am formalizing Higham's Chapter 4 Kahan summation theorem in Lean.  Please
answer as a mathematical proof adviser: give a rigorous, formalizable recurrence
or explain precisely why the proposed route cannot prove the stated bound.  Do
not rely on unstated floating-point folklore.

\subsection*{Local model}

One Kahan coefficient step for a prior source coefficient is represented by
\[
  s' = A s + B e,\qquad e' = C s + D e
\]
for homogeneous propagation, and the current input injection has coefficients
\[
  s_{\mathrm{src}} = B,\qquad e_{\mathrm{src}} = D.
\]
In paired coordinates \(T=s+e,\ E=e\), the homogeneous propagation is
\[
  T' = aT + rE,\qquad E'=cT + qE,
\]
where
\[
  a=A+C,\quad b=B+D,\quad r=(B+D)-(A+C),\quad c=C,\quad q=D-C.
\]
The returned stored-sum coefficient is \(s=T-E\).

The Lean development already proves the following local bounds for every step,
with unit roundoff \(u\) and \(0\le u\le 1\):
\[
  |a-1|\le 3u^2,\qquad |b-1|\le 2u+9u^2,
\]
\[
  |c|\le u(1+u)^2,\qquad |q|\le u(1+u)^2(3+u),
\]
\[
  |q+\delta_{\rm sub}|\le 7u^2,\qquad
  |r-q-\delta_y|\le u^2,
\]
and therefore
\[
  |r+\delta_{\rm sub}-\delta_y|\le 8u^2.
\]
There is also an exact-subtraction special case \(\delta_{\rm sub}=0\) that
already gives the desired source-shaped returned coefficient bound.  That is
not enough for arbitrary Kahan summation.

\subsection*{Known obstruction}

GPT-5.5 Pro previously suggested, and Lean now proves, a concrete p=5
round-to-even trace for inputs \((-22,-68,-248,8)\) returning \(s=-352\) while
the exact sum is \(-330\).  Hence any source-weight representation needs
\[
  \max_i |\mu_i| \ge \frac{11}{173} > \frac1{16}=2u.
\]
Lean also proves the conditional result that any same-identity trace at scale
\(p>26\) cannot satisfy a pure first-order \(2/p\) coefficient cap.  Therefore
the correct target must include a genuine \(O(nu^2)\) term; do not try to prove
a pure \(2u\) theorem.

\subsection*{Target proof shape}

For an input injected at step \(i\), define its paired source coefficient after
the own step as
\[
  T_i=b_i,\qquad E_i=D_i,
\]
then propagate through suffix steps \(j=i+1,\dots,n\) by
\[
  T_{j}=a_jT_{j-1}+r_jE_{j-1},\qquad
  E_j=c_jT_{j-1}+q_jE_{j-1}.
\]
The returned source coefficient is \(S_j=T_j-E_j\).

The existing Lean bridge can finish Higham's equation (4.8) if we can prove a
bound of the form
\[
  |S_n-1| \le 2u + C(n-i)u^2
\]
under explicit smallness assumptions such as \(nu\le 1/10\) or \(nu^2\le 1/2\),
with a concrete loose constant \(C\).  A slightly weaker but source-shaped
\(2u+Cnu^2\) theorem is fine.

\subsection*{Question}

Please provide one of the following:

1. A detailed recurrence proof that bounds \( |S_n-1| \le 2u+C(n-i)u^2 \)
   from the local facts above, including the invariant to track and every
   algebraic cancellation needed.  Constants may be loose but must be explicit.
2. A proof that these local facts are insufficient, with a small symbolic
   counterexample recurrence showing why a first-order term remains.
3. A corrected theorem with the missing hidden hypothesis stated explicitly.

The answer should be usable to write Lean lemmas.  Please phrase inequalities
in elementary real arithmetic, avoid asymptotic notation in the final bound,
and identify the exact recurrence invariant.

\end{document}
